\question
Один европейский астроном, живший в 19 веке, обнаружил новое созвездие, представляющее собой граф-дерево. Он захотел как можно быстрее поделиться своим открытием с коллегой, но была проблема: коллега жил в Америке, и единственным доступным средством связи был телеграф. Для передачи информации он решил закодировать структуру созвездия в бинарный код. Вы являетесь тем самым коллегой и должны раскодировать полученное сообщение и нарисовать полученное созвездие: ($111110101010000011001111001000$)

а) Вы захотели занести полученное созвездие в справочник, и вам необходимо составить его формальное описание в терминах графов. Посчитайте для полученного дерева:
\begin{enumerate}
\item радиус;
\item диаметр;
\item центр (отметьте на дереве);
\item число листьев (отметьте на дереве);
\end{enumerate}

б) Посмотрев в телескоп и отыскав на небе полученное созвездие, вы поняли, что ваш европейский коллега недоглядел несколько звёзд, и как следствие построил граф созвездия не до конца. Дополните дерево 6 (шестью) рёбрами так, чтобы число листьев и центров не изменилось:
\begin{enumerate}
\item нарисуйте новое дерево отдельно (покажите, что условия выполнены);
\item посчитайте для нового дерева: радиус, диаметр и центр (отметьте на дереве);
\item отметьте все листья на новом дереве;
\item закодируйте новое дерево;
\end{enumerate}

в) Обновив информацию о полученном созвездии и добавив его в справочник, вы захотели посмотреть на его связь с другими звёздами, для удобства пользуясь опять же графами. На основе уже имеющихся созвездий (деревьев) постройте следующие графы:
\begin{enumerate}
\item граф, который имеет 5 циклов разной длины
\item граф, который имеет 9 компонент вершинной двусвязности;
\item граф, который имеет 11 компонент реберной двусвязности;
\item граф, у котрого цикломатическое число 7;
\end{enumerate}

---------------

Автор -- Тимур Гонтарь
